\documentclass[12pt,a4paper]{article}

\usepackage[T2A]{fontenc}
\usepackage[utf8]{inputenc}
\usepackage[russian]{babel}
\usepackage{amsmath,amssymb,amsthm}
\usepackage{array}
\usepackage{booktabs}
\usepackage{enumitem}
\usepackage{listings}
\usepackage{xcolor}
\usepackage{geometry}
\usepackage{hyperref}

\geometry{left=25mm,right=15mm,top=20mm,bottom=20mm}

\newtheorem{definition}{Определение}
\newtheorem{theorem}{Теорема}
\newtheorem{lemma}{Лемма}
\newtheorem{remark}{Замечание}

\lstdefinestyle{code}{
    basicstyle=\ttfamily\small,
    keywordstyle=\color{blue},
    commentstyle=\color{gray},
    stringstyle=\color{red!60!black},
    breaklines=true,
    frame=single
}

\title{Билет №69 \\ \small Оптимизация многозадачной работы компьютеров. Операционные системы Windows, Unix, Linux. Особенности организации, предоставляемые услуги пользовательского взаимодействия. (ИТМО)}
\author{Сериков Владислав Александрович}
\date{16 мая 2026}

\begin{document}

\maketitle
\tableofcontents
\newpage

\section{Понятие операционной системы}

\begin{definition}
\textbf{Операционная система} --- это комплекс системных программ, управляющих аппаратными ресурсами компьютера и предоставляющих прикладным программам и пользователям удобные абстракции для работы с вычислительной системой.
\end{definition}

Операционная система выполняет две основные функции:

\begin{enumerate}
    \item \textbf{Расширенная машина}: скрывает детали аппаратуры, предоставляя удобные абстракции: процесс, файл, каталог, виртуальная память, сокет, поток, устройство.
    \item \textbf{Менеджер ресурсов}: распределяет процессорное время, оперативную память, дисковое пространство, устройства ввода-вывода, сетевые ресурсы между программами и пользователями.
\end{enumerate}

К основным задачам ОС относятся:

\begin{itemize}
    \item управление процессами и потоками;
    \item планирование выполнения задач;
    \item управление оперативной и виртуальной памятью;
    \item управление файловыми системами;
    \item управление устройствами ввода-вывода;
    \item обеспечение защиты и разграничения доступа;
    \item поддержка пользовательского интерфейса;
    \item сетевое взаимодействие;
    \item обработка ошибок и ведение журналов.
\end{itemize}

\section{Многозадачность}

\subsection{Основные определения}

\begin{definition}
\textbf{Многозадачность} --- это способность операционной системы обеспечивать видимое или фактическое одновременное выполнение нескольких задач.
\end{definition}

Если процессор один, то задачи выполняются не физически одновременно, а попеременно: ОС быстро переключает процессор между процессами. У пользователя создаётся впечатление параллельной работы.

Если процессоров или ядер несколько, то возможен настоящий параллелизм: несколько потоков действительно выполняются одновременно.

\begin{definition}
\textbf{Процесс} --- это экземпляр выполняющейся программы вместе с её адресным пространством, открытыми файлами, регистрами процессора, состоянием и прочими ресурсами.
\end{definition}

\begin{definition}
\textbf{Поток} --- это единица выполнения внутри процесса. Потоки одного процесса обычно разделяют общее адресное пространство, но имеют собственный стек, регистры и состояние выполнения.
\end{definition}

\begin{definition}
\textbf{Контекст процесса} --- это совокупность данных, необходимых для приостановки и последующего продолжения выполнения процесса: значения регистров, счётчик команд, указатель стека, состояние памяти и системные структуры ядра.
\end{definition}

\begin{definition}
\textbf{Переключение контекста} --- это операция сохранения состояния текущего процесса или потока и восстановления состояния другого процесса или потока.
\end{definition}

Переключение контекста необходимо для многозадачности, но оно не является бесплатным: во время переключения полезная пользовательская программа не выполняется. Поэтому оптимизация многозадачной работы компьютера во многом сводится к разумному выбору частоты переключений и политик планирования.

\subsection{Виды многозадачности}

Различают два основных вида многозадачности:

\begin{enumerate}
    \item \textbf{Кооперативная многозадачность}.

    Процесс сам добровольно передаёт управление операционной системе. Если программа зависает или не отдаёт управление, вся система может стать неотзывчивой.

    \item \textbf{Вытесняющая многозадачность}.

    ОС может принудительно приостановить выполняющийся процесс по таймеру или по более приоритетному событию. Современные ОС, включая Windows, Unix и Linux, используют вытесняющую многозадачность.
\end{enumerate}

\subsection{Цели оптимизации многозадачной работы}

Оптимизация многозадачной работы направлена на достижение следующих целей:

\begin{itemize}
    \item высокая загрузка процессора;
    \item малое среднее время ожидания процессов;
    \item малое среднее время отклика интерактивных программ;
    \item высокая пропускная способность системы;
    \item справедливое распределение ресурсов;
    \item предотвращение голодания процессов;
    \item эффективная работа с устройствами ввода-вывода;
    \item поддержка задач реального времени;
    \item масштабируемость на многоядерных и многопроцессорных системах.
\end{itemize}

\section{Критерии эффективности планирования}

Пусть имеется набор процессов $P_1, P_2, \dots, P_n$.

Для каждого процесса можно определить:

\begin{itemize}
    \item $a_i$ --- момент поступления процесса $P_i$ в систему;
    \item $b_i$ --- длительность процессорного обслуживания;
    \item $c_i$ --- момент завершения процесса;
    \item $w_i$ --- время ожидания;
    \item $t_i$ --- полное время пребывания процесса в системе.
\end{itemize}

\begin{definition}
\textbf{Время выполнения} или \textbf{CPU burst} --- это время, в течение которого процессу требуется процессор до следующей операции ввода-вывода или завершения.
\end{definition}

\begin{definition}
\textbf{Время ожидания} процесса:
\[
w_i = t_i - b_i,
\]
где $t_i$ --- полное время пребывания процесса в системе, $b_i$ --- суммарное процессорное время процесса.
\end{definition}

\begin{definition}
\textbf{Время оборота} или \textbf{turnaround time}:
\[
t_i = c_i - a_i.
\]
\end{definition}

\begin{definition}
\textbf{Время отклика} --- это время от поступления запроса до первого получения процессорного времени или первого видимого ответа системы.
\end{definition}

Основные критерии:

\begin{enumerate}
    \item \textbf{CPU utilization} --- процент времени, когда процессор занят полезной работой.
    \item \textbf{Throughput} --- число завершённых процессов в единицу времени.
    \item \textbf{Turnaround time} --- среднее время полного прохождения процесса через систему.
    \item \textbf{Waiting time} --- среднее время ожидания в очереди готовых процессов.
    \item \textbf{Response time} --- время реакции системы на запрос пользователя.
    \item \textbf{Fairness} --- справедливость распределения процессорного времени.
\end{enumerate}

\section{Модель процессов и состояний}

Процесс в ОС обычно проходит через следующие состояния:

\begin{enumerate}
    \item \textbf{New} --- процесс создаётся.
    \item \textbf{Ready} --- процесс готов к выполнению и ожидает процессор.
    \item \textbf{Running} --- процесс выполняется на процессоре.
    \item \textbf{Waiting / Blocked} --- процесс ожидает событие, например завершение ввода-вывода.
    \item \textbf{Terminated} --- процесс завершён.
\end{enumerate}

Типичные переходы:

\[
\text{New} \rightarrow \text{Ready} \rightarrow \text{Running} \rightarrow \text{Terminated}.
\]

Дополнительно:

\[
\text{Running} \rightarrow \text{Waiting}
\]

при запросе ввода-вывода;

\[
\text{Waiting} \rightarrow \text{Ready}
\]

при завершении ожидаемого события;

\[
\text{Running} \rightarrow \text{Ready}
\]

при вытеснении процессом с более высоким приоритетом или по истечении кванта времени.

\section{Планирование процессов}

\subsection{Уровни планирования}

В классической модели выделяют три уровня планирования:

\begin{enumerate}
    \item \textbf{Долгосрочное планирование} --- решает, какие задания допустить в систему.
    \item \textbf{Среднесрочное планирование} --- может временно выгружать процессы из памяти на диск и возвращать их обратно.
    \item \textbf{Краткосрочное планирование} --- выбирает следующий готовый процесс для выполнения на процессоре.
\end{enumerate}

В современных ОС основное значение имеет краткосрочное планирование.

\subsection{Планировщик и диспетчер}

\begin{definition}
\textbf{Планировщик} --- компонент ОС, принимающий решение о том, какой процесс или поток должен выполняться следующим.
\end{definition}

\begin{definition}
\textbf{Диспетчер} --- компонент ОС, который фактически передаёт управление выбранному процессу: выполняет переключение контекста, переводит процессор в пользовательский режим и возобновляет выполнение.
\end{definition}

Время, затрачиваемое диспетчером на передачу управления, называется \textbf{диспетчерской задержкой}. Чем она меньше, тем эффективнее система.

\section{Классические алгоритмы планирования}

\subsection{FCFS: First-Come, First-Served}

\begin{definition}
\textbf{FCFS} --- алгоритм планирования, при котором процессы выполняются в порядке поступления.
\end{definition}

\subsubsection*{Шаги алгоритма}

\begin{enumerate}
    \item Все готовые процессы помещаются в очередь FIFO.
    \item Когда процессор освобождается, выбирается процесс из начала очереди.
    \item Процесс выполняется до завершения или блокировки.
    \item Новые процессы добавляются в конец очереди.
\end{enumerate}

\subsubsection*{Пример}

Пусть процессы поступили одновременно:

\[
P_1: 10,\quad P_2: 2,\quad P_3: 1.
\]

Диаграмма Ганта:

\[
\begin{array}{c|c|c|c}
0 & P_1 & P_2 & P_3 \\
\hline
0 & 10 & 12 & 13
\end{array}
\]

Времена ожидания:

\[
w_1 = 0,\quad w_2 = 10,\quad w_3 = 12.
\]

Среднее время ожидания:

\[
\overline{w} = \frac{0 + 10 + 12}{3} = \frac{22}{3} \approx 7.33.
\]

Недостаток FCFS --- эффект конвоя: короткие процессы могут долго ждать завершения длинного процесса.

\subsection{SJF: Shortest Job First}

\begin{definition}
\textbf{SJF} --- алгоритм, выбирающий процесс с минимальным ожидаемым временем следующего процессорного обслуживания.
\end{definition}

Существуют две формы:

\begin{enumerate}
    \item невытесняющая SJF;
    \item вытесняющая SJF, также называемая SRTF --- Shortest Remaining Time First.
\end{enumerate}

\subsubsection*{Шаги невытесняющего SJF}

\begin{enumerate}
    \item Из очереди готовых процессов выбирается процесс с наименьшим CPU burst.
    \item Процесс выполняется до завершения или блокировки.
    \item После освобождения процессора выбор повторяется.
\end{enumerate}

\subsubsection*{Пример}

Пусть процессы поступили одновременно:

\[
P_1: 10,\quad P_2: 2,\quad P_3: 1.
\]

Порядок SJF:

\[
P_3 \rightarrow P_2 \rightarrow P_1.
\]

Диаграмма:

\[
\begin{array}{c|c|c|c}
0 & P_3 & P_2 & P_1 \\
\hline
0 & 1 & 3 & 13
\end{array}
\]

Времена ожидания:

\[
w_3 = 0,\quad w_2 = 1,\quad w_1 = 3.
\]

Среднее время ожидания:

\[
\overline{w} = \frac{0 + 1 + 3}{3} = \frac{4}{3} \approx 1.33.
\]

\subsection{Оптимальность SJF}

\begin{theorem}
Для множества процессов, одновременно доступных в момент времени $0$, при отсутствии вытеснения и одинаковых приоритетах алгоритм SJF минимизирует среднее время ожидания.
\end{theorem}

\begin{proof}
Рассмотрим произвольное расписание двух соседних процессов $A$ и $B$, выполняемых подряд. Пусть длительности их выполнения равны $a$ и $b$, причём $a > b$.

Пусть перед этими двумя процессами уже прошло время $T$. Тогда в порядке $A \rightarrow B$ вклады этих двух процессов во время ожидания равны:
\[
W_{AB} = T + (T + a) = 2T + a.
\]

Если поменять их местами, получаем порядок $B \rightarrow A$. Тогда:
\[
W_{BA} = T + (T + b) = 2T + b.
\]

Так как $a > b$, то
\[
W_{AB} = 2T + a > 2T + b = W_{BA}.
\]

Следовательно, если в расписании более длинная задача стоит перед более короткой, их перестановка уменьшает суммарное время ожидания.

Повторяя такие попарные перестановки, можно получить расписание, в котором процессы упорядочены по неубыванию длительности. Именно такое расписание строит SJF.

Так как каждая перестановка уменьшает или не увеличивает суммарное время ожидания, расписание SJF имеет минимальное суммарное, а значит и минимальное среднее время ожидания.
\end{proof}

\subsection{Round Robin}

\begin{definition}
\textbf{Round Robin} --- вытесняющий алгоритм планирования, при котором каждому процессу циклически выделяется квант процессорного времени $q$.
\end{definition}

\subsubsection*{Шаги алгоритма}

\begin{enumerate}
    \item Все готовые процессы помещаются в очередь FIFO.
    \item Первый процесс из очереди получает процессор на время не более $q$.
    \item Если процесс завершается раньше, он покидает систему.
    \item Если процесс не завершился за квант $q$, он вытесняется и помещается в конец очереди.
    \item Планировщик выбирает следующий процесс из начала очереди.
\end{enumerate}

\subsubsection*{Пример}

Пусть:

\[
P_1: 5,\quad P_2: 3,\quad P_3: 1,\quad q = 2.
\]

Выполнение:

\[
P_1(2), P_2(2), P_3(1), P_1(2), P_2(1), P_1(1).
\]

Диаграмма:

\[
\begin{array}{c|c|c|c|c|c|c}
0 & P_1 & P_2 & P_3 & P_1 & P_2 & P_1 \\
\hline
0 & 2 & 4 & 5 & 7 & 8 & 9
\end{array}
\]

Round Robin хорошо подходит для интерактивных систем. Если квант слишком мал, возрастает число переключений контекста. Если квант слишком велик, алгоритм приближается к FCFS.

\subsection{Планирование по приоритетам}

\begin{definition}
\textbf{Приоритетное планирование} --- алгоритм, при котором каждому процессу назначается приоритет, и процессор получает процесс с наивысшим приоритетом.
\end{definition}

Приоритеты бывают:

\begin{itemize}
    \item статические;
    \item динамические;
    \item пользовательские;
    \item системные;
    \item абсолютные;
    \item относительные.
\end{itemize}

Главная проблема --- \textbf{голодание}: процесс с низким приоритетом может очень долго не получать процессор.

Для борьбы с голоданием используется \textbf{старение}.

\begin{definition}
\textbf{Старение} --- это механизм постепенного повышения приоритета процесса по мере увеличения времени его ожидания.
\end{definition}

\subsection{Многоуровневые очереди}

При многоуровневом планировании процессы распределяются по нескольким очередям:

\begin{itemize}
    \item системные процессы;
    \item интерактивные процессы;
    \item пакетные процессы;
    \item фоновые процессы.
\end{itemize}

Каждая очередь может иметь собственный алгоритм планирования.

Например:

\[
\begin{array}{c|c|c}
\text{Очередь} & \text{Тип процессов} & \text{Алгоритм} \\
\hline
Q_0 & \text{интерактивные} & \text{Round Robin, малый квант} \\
Q_1 & \text{пользовательские} & \text{Round Robin, больший квант} \\
Q_2 & \text{фоновые} & \text{FCFS}
\end{array}
\]

\subsection{Многоуровневая очередь с обратной связью}

\begin{definition}
\textbf{MLFQ} --- многоуровневая очередь с обратной связью, в которой процесс может перемещаться между очередями в зависимости от своего поведения.
\end{definition}

Идея:

\begin{itemize}
    \item интерактивные процессы, часто выполняющие ввод-вывод, остаются в верхних очередях;
    \item CPU-bound процессы постепенно опускаются в нижние очереди;
    \item старение может возвращать долго ожидающие процессы вверх.
\end{itemize}

\subsubsection*{Шаги алгоритма MLFQ}

\begin{enumerate}
    \item Новый процесс помещается в верхнюю очередь.
    \item Планировщик всегда выбирает процесс из самой высокой непустой очереди.
    \item Если процесс использовал весь квант, он понижается в очередь ниже.
    \item Если процесс заблокировался на ввод-вывод до конца кванта, он может остаться в той же очереди.
    \item Периодически все процессы могут повышаться для предотвращения голодания.
\end{enumerate}

\section{Очереди, загрузка и закон Литтла}

\subsection{Интуитивная модель}

Многозадачная система может рассматриваться как система массового обслуживания. Процессы приходят, становятся в очередь, получают обслуживание процессором или устройством ввода-вывода, затем либо завершаются, либо снова становятся в очередь.

Для анализа используются:

\begin{itemize}
    \item $\lambda$ --- интенсивность поступления заявок;
    \item $W$ --- среднее время пребывания заявки в системе;
    \item $L$ --- среднее число заявок в системе.
\end{itemize}

\begin{theorem}[Закон Литтла]
Для стабильной системы массового обслуживания среднее число заявок в системе равно произведению интенсивности поступления заявок на среднее время пребывания заявки в системе:
\[
L = \lambda W.
\]
\end{theorem}

\begin{proof}
Рассмотрим большой интервал времени $[0,T]$.

Пусть за это время в систему поступило $N(T)$ заявок. Обозначим через $W_i$ время пребывания $i$-й заявки в системе.

Суммарное время, проведённое всеми заявками в системе, равно:
\[
\sum_{i=1}^{N(T)} W_i.
\]

С другой стороны, если $L(t)$ --- число заявок в системе в момент времени $t$, то площадь под графиком $L(t)$ на интервале $[0,T]$ равна той же величине:
\[
\int_0^T L(t)\,dt = \sum_{i=1}^{N(T)} W_i,
\]
с точностью до граничных эффектов для заявок, пришедших до $0$ или не завершившихся к $T$. В стабильной системе при $T \to \infty$ вклад граничных эффектов становится пренебрежимо малым.

Разделим обе части на $T$:
\[
\frac{1}{T}\int_0^T L(t)\,dt =
\frac{N(T)}{T} \cdot
\frac{1}{N(T)}\sum_{i=1}^{N(T)} W_i.
\]

При $T \to \infty$ левая часть стремится к среднему числу заявок в системе $L$:
\[
\frac{1}{T}\int_0^T L(t)\,dt \to L.
\]

Первый множитель справа стремится к интенсивности поступления:
\[
\frac{N(T)}{T} \to \lambda.
\]

Второй множитель справа стремится к среднему времени пребывания заявки:
\[
\frac{1}{N(T)}\sum_{i=1}^{N(T)} W_i \to W.
\]

Следовательно:
\[
L = \lambda W.
\]
\end{proof}

\subsection{Применение закона Литтла к ОС}

Если в системе в среднем находится $L = 20$ процессов, а среднее время пребывания процесса равно $W = 10$ секунд, то интенсивность завершения процессов:

\[
\lambda = \frac{L}{W} = \frac{20}{10} = 2
\]

процесса в секунду.

Закон Литтла показывает, что при увеличении числа процессов без роста пропускной способности системы увеличивается среднее время ожидания.

\section{Оптимизация многозадачности}

\subsection{Баланс CPU-bound и I/O-bound процессов}

Процессы можно условно разделить на два типа:

\begin{enumerate}
    \item \textbf{CPU-bound} --- большую часть времени используют процессор.
    \item \textbf{I/O-bound} --- часто ожидают ввод-вывод.
\end{enumerate}

Для эффективной многозадачности важно чередовать эти типы процессов. Пока один процесс ждёт диск или сеть, другой может использовать процессор.

\subsection{Оптимальный размер кванта}

Квант времени должен быть:

\begin{itemize}
    \item достаточно малым, чтобы обеспечить хорошее время отклика;
    \item достаточно большим, чтобы не тратить слишком много времени на переключения контекста.
\end{itemize}

Пусть:

\begin{itemize}
    \item $q$ --- квант времени;
    \item $s$ --- время переключения контекста.
\end{itemize}

Тогда грубая доля накладных расходов:

\[
\frac{s}{q+s}.
\]

Например, если $q = 10$ мс, $s = 0.1$ мс, то:

\[
\frac{0.1}{10 + 0.1} \approx 0.0099 \approx 1\%.
\]

Если же $q = 1$ мс, то:

\[
\frac{0.1}{1 + 0.1} \approx 0.091 \approx 9.1\%.
\]

Следовательно, слишком малый квант может сильно снизить производительность.

\subsection{Снижение числа переключений контекста}

Для уменьшения накладных расходов используются:

\begin{itemize}
    \item адаптивный выбор кванта;
    \item группировка связанных потоков;
    \item привязка потоков к ядрам;
    \item уменьшение числа активных фоновых процессов;
    \item асинхронный ввод-вывод;
    \item использование пулов потоков вместо частого создания новых потоков.
\end{itemize}

\subsection{Кэш-локальность и affinity}

\begin{definition}
\textbf{Процессорная привязка} или \textbf{CPU affinity} --- это стремление выполнять поток на том же процессорном ядре, где он выполнялся ранее.
\end{definition}

Это полезно, потому что данные потока могут оставаться в кэше процессора. Перенос потока на другое ядро может привести к кэш-промахам и снижению производительности.

\subsection{Балансировка нагрузки}

В многоядерной системе планировщик должен распределять потоки между ядрами так, чтобы:

\begin{itemize}
    \item не было перегруженных ядер;
    \item не было простаивающих ядер;
    \item сохранялась кэш-локальность;
    \item учитывалась неоднородность процессоров, если она есть.
\end{itemize}

Балансировка может быть:

\begin{itemize}
    \item \textbf{push} --- перегруженное ядро передаёт задачи другим ядрам;
    \item \textbf{pull} --- свободное ядро забирает задачи у перегруженных ядер.
\end{itemize}

\subsection{Асинхронный ввод-вывод}

Асинхронный ввод-вывод позволяет процессу не блокироваться при ожидании завершения операции. Вместо этого процесс может продолжить работу, а ОС уведомит его о завершении операции.

Это особенно важно для серверных приложений, где тысячи соединений не должны требовать тысячи постоянно активных потоков.

\subsection{Пулы потоков}

\begin{definition}
\textbf{Пул потоков} --- это заранее созданный набор рабочих потоков, которым передаются задачи на выполнение.
\end{definition}

Преимущества:

\begin{itemize}
    \item уменьшаются расходы на создание и уничтожение потоков;
    \item ограничивается максимальное число одновременно работающих потоков;
    \item улучшается управляемость нагрузки;
    \item повышается стабильность серверных приложений.
\end{itemize}

\section{Синхронизация в многозадачных ОС}

\subsection{Проблема гонки}

\begin{definition}
\textbf{Состояние гонки} --- ситуация, когда результат выполнения программы зависит от относительного порядка выполнения потоков.
\end{definition}

Пример:

\begin{lstlisting}[style=code]
x = x + 1;
\end{lstlisting}

На уровне машинных команд это может быть:

\begin{enumerate}
    \item загрузить $x$ в регистр;
    \item увеличить регистр;
    \item записать регистр обратно в $x$.
\end{enumerate}

Если два потока выполняют эти действия одновременно, одно из увеличений может потеряться.

\subsection{Критическая секция}

\begin{definition}
\textbf{Критическая секция} --- участок программы, где происходит доступ к разделяемому ресурсу.
\end{definition}

Требования к решению задачи критической секции:

\begin{enumerate}
    \item \textbf{Взаимное исключение}: одновременно в критической секции может быть не более одного процесса.
    \item \textbf{Прогресс}: если критическая секция свободна, выбор следующего входящего не должен бесконечно откладываться.
    \item \textbf{Ограниченное ожидание}: процесс не должен ждать входа в критическую секцию бесконечно долго.
\end{enumerate}

\subsection{Мьютексы, семафоры, мониторы}

\begin{definition}
\textbf{Мьютекс} --- объект синхронизации, обеспечивающий взаимное исключение.
\end{definition}

\begin{definition}
\textbf{Семафор} --- целочисленный объект синхронизации с двумя атомарными операциями: $P$ и $V$, или wait и signal.
\end{definition}

Операция wait:

\[
P(S):
\begin{cases}
S := S - 1, & \text{если } S > 0;\\
\text{блокировка процесса}, & \text{если } S = 0.
\end{cases}
\]

Операция signal:

\[
V(S): S := S + 1
\]

и, если есть ожидающие процессы, один из них пробуждается.

\begin{definition}
\textbf{Монитор} --- высокоуровневая конструкция синхронизации, объединяющая данные, процедуры работы с ними и механизм взаимного исключения.
\end{definition}

\section{Взаимные блокировки}

\begin{definition}
\textbf{Взаимная блокировка} или \textbf{deadlock} --- ситуация, когда множество процессов ожидает события, которое может быть вызвано только другим процессом из этого же множества.
\end{definition}

\subsection{Условия Коффмана}

Взаимная блокировка возможна при одновременном выполнении четырёх условий:

\begin{enumerate}
    \item \textbf{Взаимное исключение}: некоторый ресурс не может использоваться более чем одним процессом одновременно.
    \item \textbf{Удержание и ожидание}: процесс удерживает одни ресурсы и ожидает другие.
    \item \textbf{Отсутствие принудительного отъёма}: ресурс нельзя отобрать у процесса насильно.
    \item \textbf{Циклическое ожидание}: существует цикл процессов, каждый из которых ожидает ресурс, удерживаемый следующим процессом в цикле.
\end{enumerate}

\subsection{Методы борьбы с deadlock}

\begin{enumerate}
    \item \textbf{Предотвращение} --- нарушение хотя бы одного условия Коффмана.
    \item \textbf{Избежание} --- анализ будущих состояний, например алгоритм банкира.
    \item \textbf{Обнаружение и восстановление} --- периодический поиск deadlock и принудительное завершение или откат процессов.
    \item \textbf{Игнорирование} --- если deadlock крайне редок, ОС может не применять дорогостоящие механизмы предотвращения.
\end{enumerate}

\section{Управление памятью}

\subsection{Виртуальная память}

\begin{definition}
\textbf{Виртуальная память} --- механизм, при котором каждому процессу предоставляется собственное логическое адресное пространство, отображаемое ОС и аппаратурой на физическую память и внешнее хранилище.
\end{definition}

Преимущества виртуальной памяти:

\begin{itemize}
    \item изоляция процессов;
    \item возможность запускать программы, превышающие размер физической памяти;
    \item упрощение программирования;
    \item защита памяти;
    \item совместное использование библиотек;
    \item эффективная загрузка страниц по требованию.
\end{itemize}

\subsection{Страничная организация}

Память делится на:

\begin{itemize}
    \item \textbf{страницы} --- блоки виртуальной памяти;
    \item \textbf{кадры страниц} --- блоки физической памяти.
\end{itemize}

Отображение задаётся таблицами страниц.

Виртуальный адрес можно представить как:

\[
\text{виртуальный адрес} = (\text{номер страницы}, \text{смещение}).
\]

Физический адрес:

\[
\text{физический адрес} = (\text{номер кадра}, \text{смещение}).
\]

\subsection{TLB}

\begin{definition}
\textbf{TLB} --- буфер трансляции адресов, кэш недавних преобразований виртуальных адресов в физические.
\end{definition}

TLB уменьшает накладные расходы на обращение к таблицам страниц.

\subsection{Замещение страниц}

Если нужной страницы нет в физической памяти, возникает \textbf{page fault}. ОС должна загрузить страницу с диска, возможно вытеснив другую страницу.

Алгоритмы замещения:

\begin{itemize}
    \item FIFO;
    \item LRU;
    \item Clock;
    \item оптимальный алгоритм Белади.
\end{itemize}

\subsection{Оптимальный алгоритм замещения страниц}

\begin{definition}
\textbf{Оптимальный алгоритм замещения страниц} вытесняет ту страницу, которая не понадобится дольше всего в будущем.
\end{definition}

Этот алгоритм практически нереализуем в общем случае, так как требует знания будущего, но используется как теоретический эталон.

\begin{theorem}
Оптимальный алгоритм замещения страниц минимизирует число страничных промахов среди всех алгоритмов с тем же числом кадров памяти.
\end{theorem}

\begin{proof}
Рассмотрим произвольный алгоритм $A$ и оптимальный алгоритм $OPT$. Пусть они впервые принимают разные решения при вытеснении страницы.

Алгоритм $OPT$ вытесняет страницу $p$, которая будет использована позже всех остальных страниц, находящихся в памяти. Алгоритм $A$ вытесняет другую страницу $q$.

Так как $p$ будет использована не раньше, чем $q$, до момента следующего обращения к $p$ алгоритм, который вытеснил $p$, не хуже алгоритма, который вытеснил $q$. Если до обращения к $p$ потребуется загрузить новые страницы, можно построить расписание замен, совпадающее с решениями $A$, но не дающее большего числа промахов.

Таким образом, любое первое отклонение от выбора $OPT$ можно заменить выбором $OPT$ без увеличения числа промахов. Повторяя этот аргумент для всех отличающихся решений, получаем, что существует оптимальное расписание замен, совпадающее с алгоритмом, вытесняющим страницу, которая понадобится позже всех.

Следовательно, такой алгоритм минимизирует число страничных промахов.
\end{proof}

\section{Файловые системы}

\subsection{Файл и каталог}

\begin{definition}
\textbf{Файл} --- именованная совокупность данных на внешнем носителе.
\end{definition}

\begin{definition}
\textbf{Каталог} --- специальная структура файловой системы, связывающая имена файлов с их метаданными и расположением.
\end{definition}

ОС предоставляет операции:

\begin{itemize}
    \item создание файла;
    \item удаление файла;
    \item открытие;
    \item закрытие;
    \item чтение;
    \item запись;
    \item перемещение указателя;
    \item изменение прав доступа.
\end{itemize}

\subsection{Метаданные}

Метаданные файла могут включать:

\begin{itemize}
    \item имя;
    \item тип;
    \item размер;
    \item владельца;
    \item права доступа;
    \item временные метки;
    \item адреса блоков на диске.
\end{itemize}

\subsection{Журналирование}

\begin{definition}
\textbf{Журналируемая файловая система} --- файловая система, которая перед изменением основных структур записывает намерение или описание операции в журнал.
\end{definition}

Журналирование повышает устойчивость к сбоям: после аварийного выключения ОС может восстановить согласованность файловой системы, просмотрев журнал.

\section{Операционная система Unix}

\subsection{Общая характеристика}

Unix --- семейство многопользовательских, многозадачных операционных систем, оказавшее огромное влияние на развитие современных ОС.

Ключевые идеи Unix:

\begin{itemize}
    \item всё представляется как файл;
    \item небольшие программы должны хорошо выполнять одну задачу;
    \item программы можно соединять конвейерами;
    \item текстовые интерфейсы и файлы конфигурации;
    \item иерархическая файловая система;
    \item строгая модель пользователей и прав доступа;
    \item переносимость за счёт реализации на языке C.
\end{itemize}

\subsection{Архитектура Unix}

Упрощённо Unix можно представить в виде уровней:

\[
\begin{array}{c}
\text{Пользовательские программы} \\
\hline
\text{Командная оболочка и библиотеки} \\
\hline
\text{Системные вызовы} \\
\hline
\text{Ядро} \\
\hline
\text{Аппаратное обеспечение}
\end{array}
\]

Ядро Unix отвечает за:

\begin{itemize}
    \item процессы;
    \item память;
    \item файловые системы;
    \item устройства;
    \item межпроцессное взаимодействие;
    \item сетевые протоколы.
\end{itemize}

\subsection{Процессы в Unix}

В Unix процесс создаётся с помощью системного вызова:

\begin{lstlisting}[style=code]
fork()
\end{lstlisting}

Он создаёт копию текущего процесса. После этого часто вызывается:

\begin{lstlisting}[style=code]
exec()
\end{lstlisting}

для замены образа процесса новой программой.

Ожидание завершения дочернего процесса выполняется через:

\begin{lstlisting}[style=code]
wait()
\end{lstlisting}

\subsubsection*{Минимальный пример}

\begin{lstlisting}[style=code,language=C]
#include <unistd.h>
#include <sys/wait.h>
#include <stdio.h>

int main() {
    pid_t pid = fork();

    if (pid == 0) {
        execlp("ls", "ls", "-l", NULL);
    } else {
        wait(NULL);
        printf("Child finished\n");
    }

    return 0;
}
\end{lstlisting}

Смысл:

\begin{enumerate}
    \item родительский процесс вызывает \texttt{fork()};
    \item появляется дочерний процесс;
    \item дочерний процесс заменяет себя программой \texttt{ls};
    \item родитель ждёт завершения дочернего процесса.
\end{enumerate}

\subsection{Файловая модель Unix}

Unix использует единую иерархическую файловую систему с корнем:

\[
/
\]

Типичные каталоги:

\[
\begin{array}{c|l}
\text{Каталог} & \text{Назначение} \\
\hline
/bin & \text{основные пользовательские команды} \\
/sbin & \text{системные команды администратора} \\
/etc & \text{конфигурационные файлы} \\
/home & \text{домашние каталоги пользователей} \\
/var & \text{изменяемые данные: журналы, очереди} \\
/tmp & \text{временные файлы} \\
/usr & \text{пользовательские программы и библиотеки} \\
/dev & \text{файлы устройств}
\end{array}
\]

\subsection{Права доступа Unix}

Для файла задаются права:

\[
r, w, x
\]

для трёх категорий:

\begin{itemize}
    \item владелец;
    \item группа;
    \item остальные.
\end{itemize}

Например:

\[
-rwxr-xr--
\]

означает:

\[
\begin{array}{c|c}
\text{Категория} & \text{Права} \\
\hline
\text{Владелец} & rwx \\
\text{Группа} & r-x \\
\text{Остальные} & r--
\end{array}
\]

Числовая запись:

\[
r = 4,\quad w = 2,\quad x = 1.
\]

Например:

\[
chmod\ 755\ file
\]

задаёт права:

\[
rwxr-xr-x.
\]

\section{Операционная система Linux}

\subsection{Общая характеристика}

Linux --- Unix-подобная операционная система, основанная на ядре Linux и обычно распространяемая в составе дистрибутивов, включающих системные библиотеки, утилиты, оболочки, графические среды и пакетные менеджеры.

Примеры дистрибутивов:

\begin{itemize}
    \item Debian;
    \item Ubuntu;
    \item Fedora;
    \item Arch Linux;
    \item openSUSE;
    \item Red Hat Enterprise Linux;
    \item AlmaLinux;
    \item Rocky Linux.
\end{itemize}

\subsection{Архитектура Linux}

Linux имеет монолитное модульное ядро. Это означает, что основные подсистемы работают в пространстве ядра, но многие драйверы и компоненты могут загружаться как модули.

Основные подсистемы ядра Linux:

\begin{itemize}
    \item планировщик процессов;
    \item управление памятью;
    \item виртуальная файловая система VFS;
    \item сетевой стек;
    \item драйверы устройств;
    \item подсистема блочного ввода-вывода;
    \item механизмы безопасности;
    \item межпроцессное взаимодействие.
\end{itemize}

\subsection{Процессы и потоки в Linux}

В Linux процессы и потоки представлены общей структурой задачи. Поток можно понимать как процесс, разделяющий часть ресурсов с другими потоками.

Создание процессов и потоков связано с системными вызовами семейства:

\begin{lstlisting}[style=code]
fork()
clone()
execve()
\end{lstlisting}

\subsection{Планирование в Linux}

Современный Linux использует сложные планировщики, поддерживающие:

\begin{itemize}
    \item обычные процессы;
    \item задачи реального времени;
    \item deadline-задачи;
    \item балансировку нагрузки по ядрам;
    \item NUMA-архитектуры;
    \item группы управления ресурсами.
\end{itemize}

Для обычных пользовательских задач долгое время ключевой идеей был справедливый планировщик, стремящийся распределять процессорное время пропорционально весам задач.

\subsection{Виртуальное время выполнения}

Идея справедливого планирования состоит в том, что у каждой задачи ведётся виртуальное время выполнения:

\[
vruntime.
\]

Чем меньше $vruntime$, тем меньше задача получила процессора относительно других, поэтому тем раньше её следует выполнить.

Упрощённый алгоритм:

\begin{enumerate}
    \item Для каждой готовой задачи хранится $vruntime$.
    \item Выбирается задача с минимальным $vruntime$.
    \item Она выполняется некоторое время.
    \item Её $vruntime$ увеличивается с учётом приоритета.
    \item Задача возвращается в структуру готовых задач.
\end{enumerate}

\subsection{Файловые системы Linux}

Linux поддерживает множество файловых систем:

\begin{itemize}
    \item ext2;
    \item ext3;
    \item ext4;
    \item XFS;
    \item Btrfs;
    \item FAT;
    \item exFAT;
    \item NTFS;
    \item tmpfs;
    \item procfs;
    \item sysfs.
\end{itemize}

Особое значение имеют псевдофайловые системы:

\begin{itemize}
    \item \texttt{/proc} --- информация о процессах и состоянии ядра;
    \item \texttt{/sys} --- информация об устройствах и драйверах;
    \item \texttt{/dev} --- файлы устройств.
\end{itemize}

\subsection{Пользовательское взаимодействие в Linux}

Linux предоставляет несколько уровней взаимодействия:

\begin{enumerate}
    \item \textbf{Командная строка}: bash, zsh, fish.
    \item \textbf{Графические оболочки}: GNOME, KDE Plasma, XFCE.
    \item \textbf{Системные утилиты}: \texttt{ps}, \texttt{top}, \texttt{htop}, \texttt{systemctl}, \texttt{journalctl}.
    \item \textbf{Пакетные менеджеры}: \texttt{apt}, \texttt{dnf}, \texttt{pacman}, \texttt{zypper}.
    \item \textbf{Удалённый доступ}: SSH.
\end{enumerate}

\section{Операционная система Windows}

\subsection{Общая характеристика}

Windows --- семейство операционных систем Microsoft для персональных компьютеров, рабочих станций, серверов и встроенных систем.

Современные версии Windows основаны на архитектуре Windows NT. Это многозадачная, многопользовательская ОС с поддержкой виртуальной памяти, защиты, сетевого взаимодействия, графического интерфейса и развитой модели безопасности.

\subsection{Архитектура Windows NT}

Упрощённая архитектура Windows NT:

\[
\begin{array}{c}
\text{Пользовательские приложения} \\
\hline
\text{Подсистемы пользовательского режима} \\
\hline
\text{Системные библиотеки} \\
\hline
\text{Исполнительная система Windows Executive} \\
\hline
\text{Ядро} \\
\hline
\text{HAL} \\
\hline
\text{Аппаратное обеспечение}
\end{array}
\]

Основные компоненты:

\begin{itemize}
    \item \textbf{Windows Executive} --- набор служб ядра высокого уровня;
    \item \textbf{Kernel} --- низкоуровневое планирование, синхронизация, обработка прерываний;
    \item \textbf{HAL} --- слой абстракции оборудования;
    \item \textbf{Object Manager} --- управление объектами ОС;
    \item \textbf{Memory Manager} --- управление виртуальной памятью;
    \item \textbf{I/O Manager} --- управление вводом-выводом;
    \item \textbf{Security Reference Monitor} --- проверка доступа;
    \item \textbf{Process Manager} --- процессы и потоки;
    \item \textbf{Configuration Manager} --- реестр;
    \item \textbf{Cache Manager} --- кэширование файлового ввода-вывода.
\end{itemize}

\subsection{Процессы и потоки в Windows}

В Windows процесс является контейнером ресурсов. Выполнение осуществляется потоками.

Процесс содержит:

\begin{itemize}
    \item виртуальное адресное пространство;
    \item дескрипторы объектов;
    \item маркер безопасности;
    \item список потоков;
    \item информацию о загруженных модулях;
    \item параметры окружения.
\end{itemize}

Поток содержит:

\begin{itemize}
    \item состояние регистров;
    \item стек пользовательского режима;
    \item стек режима ядра;
    \item приоритет;
    \item состояние планирования.
\end{itemize}

\subsection{Планирование в Windows}

Windows использует вытесняющее приоритетное планирование потоков.

Основные свойства:

\begin{itemize}
    \item планируются именно потоки, а не процессы;
    \item используются уровни приоритета;
    \item есть динамическое изменение приоритетов;
    \item интерактивные задачи могут получать повышение приоритета;
    \item поддерживается многопроцессорное планирование;
    \item поддерживаются классы приоритетов процессов и приоритеты потоков.
\end{itemize}

Типичные классы приоритета процесса:

\begin{itemize}
    \item Idle;
    \item Below Normal;
    \item Normal;
    \item Above Normal;
    \item High;
    \item Realtime.
\end{itemize}

\subsection{Файловые системы Windows}

Основная файловая система Windows --- NTFS.

Особенности NTFS:

\begin{itemize}
    \item журналирование;
    \item списки контроля доступа ACL;
    \item поддержка больших файлов;
    \item атрибуты файлов;
    \item сжатие;
    \item шифрование;
    \item жёсткие ссылки;
    \item точки повторной обработки;
    \item квоты;
    \item теневые копии на уровне системы.
\end{itemize}

Также Windows поддерживает FAT32, exFAT, ReFS и другие файловые системы в зависимости от редакции и назначения системы.

\subsection{Реестр Windows}

\begin{definition}
\textbf{Реестр Windows} --- иерархическая база данных конфигурации ОС, драйверов, служб и приложений.
\end{definition}

Основные ветви:

\[
\begin{array}{c|l}
\text{Ветвь} & \text{Назначение} \\
\hline
HKEY\_LOCAL\_MACHINE & \text{конфигурация компьютера} \\
HKEY\_CURRENT\_USER & \text{настройки текущего пользователя} \\
HKEY\_CLASSES\_ROOT & \text{ассоциации файлов и COM} \\
HKEY\_USERS & \text{профили пользователей} \\
HKEY\_CURRENT\_CONFIG & \text{текущая аппаратная конфигурация}
\end{array}
\]

\subsection{Пользовательское взаимодействие в Windows}

Windows предоставляет развитую систему пользовательского взаимодействия:

\begin{enumerate}
    \item \textbf{Графический интерфейс}: рабочий стол, окна, меню, панель задач, проводник.
    \item \textbf{Командная строка}: \texttt{cmd.exe}.
    \item \textbf{PowerShell}: объектно-ориентированная командная оболочка и язык автоматизации.
    \item \textbf{Панель управления и приложение ``Параметры''}.
    \item \textbf{Диспетчер задач}.
    \item \textbf{Просмотр событий}.
    \item \textbf{Удалённый рабочий стол}.
    \item \textbf{Windows Terminal}.
    \item \textbf{Windows Subsystem for Linux}.
\end{enumerate}

\section{Сравнение Windows, Unix и Linux}

\subsection{Общее сравнение}

\[
\begin{array}{p{3cm}|p{4cm}|p{4cm}|p{4cm}}
\toprule
\textbf{Критерий} & \textbf{Windows} & \textbf{Unix} & \textbf{Linux} \\
\midrule
Происхождение & Microsoft, линия NT & Историческое семейство ОС & Unix-подобная ОС с ядром Linux \\
\midrule
Модель распространения & Обычно проприетарная & Разная, часто коммерческая & В основном свободная и открытая \\
\midrule
Интерфейс & Сильный акцент на GUI & CLI и серверная работа & CLI, GUI, серверы, встраиваемые системы \\
\midrule
Файловая система & Диски с буквами, NTFS & Единое дерево от / & Единое дерево от / \\
\midrule
Конфигурация & Реестр, GUI, политики, файлы & Текстовые файлы & Текстовые файлы, systemd, пакетные системы \\
\midrule
Безопасность & ACL, домены, UAC, SID & Пользователь-группа-остальные, ACL & Права Unix, ACL, capabilities, SELinux/AppArmor \\
\midrule
Планирование & Приоритетное вытесняющее & Зависит от реализации & Вытесняющее, справедливое, RT/deadline-классы \\
\midrule
Основное применение & ПК, офис, игры, корпоративная среда & Серверы, рабочие станции, промышленные системы & Серверы, облака, суперкомпьютеры, встраиваемые системы, ПК \\
\bottomrule
\end{array}
\]

\subsection{Философия организации}

\subsubsection*{Unix}

Unix ориентирован на простоту абстракций:

\begin{itemize}
    \item файл как универсальный интерфейс;
    \item конвейеры;
    \item текстовые потоки;
    \item малые утилиты;
    \item композиция программ.
\end{itemize}

\subsubsection*{Linux}

Linux наследует Unix-философию, но расширяет её:

\begin{itemize}
    \item открытая разработка;
    \item модульность ядра;
    \item огромное число поддерживаемых архитектур;
    \item контейнеризация;
    \item развитые сетевые возможности;
    \item мощные средства администрирования.
\end{itemize}

\subsubsection*{Windows}

Windows исторически ориентирована на:

\begin{itemize}
    \item удобство массового пользователя;
    \item единый графический интерфейс;
    \item совместимость приложений;
    \item интеграцию с корпоративной инфраструктурой;
    \item развитую объектную модель ядра;
    \item поддержку большого числа пользовательских программ и драйверов.
\end{itemize}

\section{Сервисы операционных систем}

ОС предоставляет пользователям и программам следующие услуги:

\begin{enumerate}
    \item \textbf{Выполнение программ}.

    ОС загружает программу в память, создаёт процесс, выделяет ресурсы и передаёт управление.

    \item \textbf{Операции ввода-вывода}.

    Программы не работают с устройствами напрямую, а используют системные вызовы.

    \item \textbf{Работа с файлами}.

    Создание, удаление, чтение, запись, поиск, изменение атрибутов.

    \item \textbf{Коммуникация}.

    Межпроцессное взаимодействие, каналы, очереди сообщений, сокеты, разделяемая память.

    \item \textbf{Обнаружение ошибок}.

    ОС обрабатывает ошибки памяти, устройств, файловых систем, сетей, программ.

    \item \textbf{Распределение ресурсов}.

    Процессор, память, диски, устройства, сеть.

    \item \textbf{Учёт использования}.

    Журналы, статистика, квоты, мониторинг.

    \item \textbf{Защита и безопасность}.

    Аутентификация, авторизация, изоляция, аудит, шифрование.
\end{enumerate}

\section{Системные вызовы}

\begin{definition}
\textbf{Системный вызов} --- контролируемый переход из пользовательского режима в режим ядра для выполнения операции, требующей привилегий ОС.
\end{definition}

Примеры категорий системных вызовов:

\[
\begin{array}{c|l}
\text{Категория} & \text{Примеры} \\
\hline
Процессы & fork, exec, wait, CreateProcess \\
Файлы & open, read, write, close, CreateFile \\
Память & mmap, brk, VirtualAlloc \\
Устройства & ioctl, DeviceIoControl \\
Сеть & socket, bind, connect, send, recv \\
Безопасность & chmod, setuid, access checks
\end{array}
\]

\section{Пользовательские интерфейсы}

\subsection{CLI}

\begin{definition}
\textbf{CLI} --- командный интерфейс, в котором пользователь взаимодействует с системой с помощью текстовых команд.
\end{definition}

Преимущества CLI:

\begin{itemize}
    \item автоматизация;
    \item скрипты;
    \item удобство удалённого управления;
    \item высокая точность операций;
    \item малые требования к ресурсам.
\end{itemize}

Недостатки:

\begin{itemize}
    \item высокий порог входа;
    \item необходимость помнить команды;
    \item меньшая наглядность для начинающих пользователей.
\end{itemize}

\subsection{GUI}

\begin{definition}
\textbf{GUI} --- графический пользовательский интерфейс, основанный на окнах, значках, меню, указателях и визуальных элементах управления.
\end{definition}

Преимущества GUI:

\begin{itemize}
    \item наглядность;
    \item простота освоения;
    \item удобство интерактивной работы;
    \item поддержка мультимедиа;
    \item развитые средства работы с документами.
\end{itemize}

Недостатки:

\begin{itemize}
    \item сложнее автоматизировать;
    \item выше требования к ресурсам;
    \item менее удобен для массового администрирования серверов.
\end{itemize}

\subsection{API и библиотеки}

Кроме прямого пользовательского интерфейса ОС предоставляет программные интерфейсы:

\begin{itemize}
    \item POSIX API в Unix/Linux-системах;
    \item Win32 API в Windows;
    \item .NET API;
    \item системные библиотеки C;
    \item графические API;
    \item сетевые API.
\end{itemize}

\section{Безопасность и защита}

\subsection{Режимы процессора}

Обычно ОС использует минимум два режима:

\begin{enumerate}
    \item \textbf{Пользовательский режим} --- ограниченные права.
    \item \textbf{Режим ядра} --- полный доступ к аппаратуре и памяти.
\end{enumerate}

Переход в режим ядра осуществляется через:

\begin{itemize}
    \item системный вызов;
    \item прерывание;
    \item исключение;
    \item аппаратное событие.
\end{itemize}

\subsection{Аутентификация и авторизация}

\begin{definition}
\textbf{Аутентификация} --- проверка личности субъекта.
\end{definition}

\begin{definition}
\textbf{Авторизация} --- проверка прав субъекта на выполнение действия.
\end{definition}

Пример:

\begin{itemize}
    \item ввод пароля --- аутентификация;
    \item проверка права на чтение файла --- авторизация.
\end{itemize}

\subsection{Модели доступа}

Unix традиционно использует модель:

\[
\text{user/group/others}.
\]

Windows активно использует ACL:

\begin{definition}
\textbf{ACL} --- список контроля доступа, определяющий, какие субъекты какие операции могут выполнять над объектом.
\end{definition}

Linux поддерживает как классические Unix-права, так и ACL, capabilities, мандатные механизмы безопасности.

\section{Виртуализация и контейнеризация}

\subsection{Виртуализация}

\begin{definition}
\textbf{Виртуализация} --- технология создания программной модели вычислительной среды, изолированной от физического оборудования.
\end{definition}

Гипервизоры бывают:

\begin{enumerate}
    \item \textbf{Тип 1} --- работают непосредственно на аппаратуре.
    \item \textbf{Тип 2} --- работают поверх обычной ОС.
\end{enumerate}

Преимущества:

\begin{itemize}
    \item изоляция;
    \item консолидация серверов;
    \item удобство тестирования;
    \item миграция виртуальных машин;
    \item снимки состояния.
\end{itemize}

\subsection{Контейнеризация}

\begin{definition}
\textbf{Контейнеризация} --- изоляция приложений на уровне ОС с использованием общего ядра.
\end{definition}

В Linux контейнеризация основана на:

\begin{itemize}
    \item namespaces;
    \item cgroups;
    \item capabilities;
    \item union/overlay файловых системах.
\end{itemize}

Контейнеры легче виртуальных машин, но используют общее ядро ОС.

\section{Практические средства наблюдения и оптимизации}

\subsection{Windows}

Полезные средства:

\begin{itemize}
    \item Диспетчер задач;
    \item Resource Monitor;
    \item Performance Monitor;
    \item Event Viewer;
    \item PowerShell;
    \item Windows Terminal;
    \item Sysinternals Process Explorer;
    \item Sysinternals Process Monitor.
\end{itemize}

Примеры команд PowerShell:

\begin{lstlisting}[style=code]
Get-Process
Get-Service
Get-EventLog -LogName System -Newest 20
\end{lstlisting}

\subsection{Linux/Unix}

Полезные команды:

\begin{lstlisting}[style=code]
ps aux
top
htop
nice
renice
kill
free -h
vmstat
iostat
sar
df -h
du -sh
journalctl
systemctl
\end{lstlisting}

Пример изменения приоритета:

\begin{lstlisting}[style=code]
nice -n 10 ./program
renice -n 5 -p 1234
\end{lstlisting}

Чем больше значение nice, тем ниже приоритет обычного процесса.

\section{Типовые проблемы многозадачности}

\subsection{Перегрузка процессора}

Симптомы:

\begin{itemize}
    \item высокая загрузка CPU;
    \item медленный отклик;
    \item рост очереди готовых процессов;
    \item нагрев и троттлинг.
\end{itemize}

Меры:

\begin{itemize}
    \item завершить лишние процессы;
    \item снизить приоритет фоновых задач;
    \item оптимизировать алгоритмы;
    \item добавить параллелизм;
    \item увеличить число ядер;
    \item использовать профилирование.
\end{itemize}

\subsection{Недостаток памяти}

Симптомы:

\begin{itemize}
    \item активное использование swap/pagefile;
    \item частые page fault;
    \item резкое падение производительности;
    \item thrashing.
\end{itemize}

\begin{definition}
\textbf{Thrashing} --- состояние, при котором система тратит большую часть времени на обмен страницами между памятью и диском, а не на полезное выполнение программ.
\end{definition}

Меры:

\begin{itemize}
    \item закрыть лишние приложения;
    \item увеличить объём RAM;
    \item уменьшить число одновременно работающих процессов;
    \item найти утечки памяти;
    \item настроить параметры swap/pagefile;
    \item использовать более эффективные структуры данных.
\end{itemize}

\subsection{Дисковая перегрузка}

Симптомы:

\begin{itemize}
    \item высокая задержка ввода-вывода;
    \item большая очередь диска;
    \item медленная загрузка приложений;
    \item задержки при обращении к файлам.
\end{itemize}

Меры:

\begin{itemize}
    \item использовать SSD;
    \item включить кэширование;
    \item оптимизировать доступ к файлам;
    \item уменьшить случайный ввод-вывод;
    \item использовать асинхронный ввод-вывод;
    \item разделять нагрузку по дискам.
\end{itemize}

\section{Итоговые положения}

Многозадачная работа компьютера обеспечивается совместной работой планировщика процессов, механизмов виртуальной памяти, подсистемы ввода-вывода, файловой системы, средств синхронизации и механизмов безопасности.

Оптимизация многозадачности требует баланса между несколькими противоречивыми целями:

\begin{itemize}
    \item максимальная загрузка процессора;
    \item минимальное время отклика;
    \item справедливость;
    \item низкие накладные расходы;
    \item предсказуемость;
    \item масштабируемость;
    \item безопасность.
\end{itemize}

Windows, Unix и Linux решают эти задачи по-разному:

\begin{itemize}
    \item Windows делает акцент на графическом интерфейсе, совместимости, объектной модели ядра и корпоративной инфраструктуре.
    \item Unix задаёт классическую модель простых абстракций, файлового интерфейса, процессов, конвейеров и текстовых инструментов.
    \item Linux развивает Unix-подход в открытой, масштабируемой и гибкой системе, широко используемой на серверах, суперкомпьютерах, в облаках, контейнерах и встраиваемых устройствах.
\end{itemize}

Главная идея современных ОС состоит в том, чтобы предоставить пользователю удобную иллюзию одновременной и безопасной работы множества программ, эффективно используя ограниченные аппаратные ресурсы.

\end{document}
